\documentclass[11pt]{article}
\usepackage[margin=1in]{geometry}
\usepackage{amsmath,amssymb}
\usepackage{booktabs}
\usepackage{hyperref}
\usepackage{xcolor}

\newcommand{\todo}[1]{\textcolor{red}{[#1]}}
\newcommand{\registered}{\textsc{registered}}
\newcommand{\notrun}{\textsc{not yet measured}}

\title{Depth-Recycled Recurrent Diffusion for Long-Context Agent Workloads}
\author{}
\date{}

\begin{document}
\maketitle

\begin{abstract}
Recurrent language models promise constant per-token state, yet the state that gives
them that promise is also the only channel through which context survives. We study a
\emph{non-latent} recurrent diffusion model: a bidirectionally scanned RWKV-7 backbone
that denoises masked tokens in place, with no learned latent sidecar and no separate
decoder. The question is whether \emph{depth recycling} --- re-applying a tied block at
increased executed depth --- improves long-context behaviour, or merely spends compute.

This manuscript reports a program at an unusual and deliberate stage, and we state it
plainly rather than implying otherwise. We contribute (i) a \emph{correction} to the
program's own prior evidence: several earlier decode-side conclusions do not survive
re-analysis, and one class of them was traced to an end-of-text id confusion that had
silently zeroed an entire evaluation; (ii) a machine-checked architecture contract that
fixes the object under study, including what it explicitly does \emph{not} yet qualify;
(iii) a rights-cleared, hash-bound asset and evaluation stack, with an offline
tokenizer qualification that settles an open id-convention conflict from the artifacts
rather than by choice; and (iv) a registered protocol --- controlled arms, a long-context
contract, and an H100 deployment measurement --- with frozen decision rules and stated
failure interpretations. The confirmatory measurements that protocol defines are
\notrun{}: no depth-recycling quality gain, no long-context macro improvement, and no
serving goodput result is reported here, and none is claimed. We consider the corrected
evidence and the frozen protocol to be the contribution, and we mark every cell where a
number is absent.
\end{abstract}

\section{Introduction}

An autoregressive recurrent model carries a fixed-size state forward and reads the next
token from it. That state is the mechanism by which context reaches the present, and it
is also a bottleneck: whatever does not fit is gone. Bidirectional scanning and
iterative denoising both attack this from different sides. Bidirectional scanning lets
every position see left \emph{and} right context within the window. Denoising lets the
model revise rather than commit. Depth recycling --- running a tied block more times
than the backbone has layers --- is a third, cheaper lever: it adds executed depth
without adding parameters.

This paper asks whether that third lever does anything, and it is organised around the
possibility that the answer is \emph{no}. Two failure modes are live and both are
unremarkable: recycling may be an expensive way to spend compute that a wider baseline
would have spent better, and any short-context benefit may simply not transfer to long
context, where a fixed-size state is the binding constraint. We therefore commit in
advance to the comparisons that would distinguish a real architectural advantage from
those two, and to the interpretation of a null result.

\paragraph{Why this manuscript looks the way it does.} The program this paper belongs to
has completed its evidence audit, its architecture contract, its asset qualification and
its runtime calibration, and has \emph{not} run the confirmatory training and evaluation
campaign. We have written the paper to that state: the completed work is reported as
results, the protocol is reported as a registrable protocol, and the cells that the
protocol defines are left visibly empty. We think this is more useful than either a
premature results paper or a plan, and the alternative --- filling the cells with
fixture numbers --- is excluded by the program's own rules.

\section{Positioning}

Depth recycling and bidirectional recurrence over SSMs are both active areas, and three
comparisons are mandatory for any novelty claim here: DiffuMamba~\cite{diffumamba},
Triplet-Block Diffusion RWKV~\cite{tripletrwkv}, and partial-bidirectional hybrid block
diffusion~\cite{partialbidir}. We make \emph{no} firstness claim. Diffusion language
models, RWKV-family recurrent models, and depth-recurrent or looped architectures all
predate this work; the contribution we aim to defend is a controlled comparison under
equal-compute and equal-quality constraints, and the corrected evidence that determines
what such a comparison should even measure.

\paragraph{What is not claimed.} We do not claim a latent semantic sidecar, a learned
coordinator, or a macro-neuron network; introducing one to rescue this line would change
the object under study. We do not claim that a causal-path negative log-likelihood gain
proves a generation improvement, nor that equal token counts imply equal compute. We do
not claim any physical edge-device latency, energy, thermal, or hardware-efficiency
result from H100 proxies: software memory caps do not reproduce edge compute or
bandwidth. And we do not claim a global long-context superiority theorem from local
spectral observations.

\section{Correcting our own prior evidence}
\label{sec:correction}

The first contribution is negative and concerns the program's earlier conclusions.

\paragraph{A decode-side id confusion had zeroed an evaluation.} In a prior bidirectional
RWKV evaluation line, generation was truncated at token id \texttt{65530}, treated as
end-of-text. Measured against the tokenizer, \texttt{65530} decodes to \texttt{'\textbackslash
n\textbackslash n'} --- an ordinary blank line --- while the real end-of-text is
\texttt{65531}, and id \texttt{0} also decodes to that same string. Because the chat
prompt ended with a code fence, any completion that opened with a blank line was
truncated to the empty string, and pass@1 was \emph{exactly} $0.0$ on 164 of 164 tasks
across three successive training-data revisions. The invariance of that number across
three rounds was itself the tell: a decode-side string bug is unreachable from the
corpus, so no amount of training-data work could have moved it. The fix is a stop set of
$\{65531, 0\}$ (plus \texttt{eos\_token\_id} when it is an integer), never \texttt{65530}.

\paragraph{An id-convention conflict, settled from the artifacts.} A separate internal
record carried an unresolved conflict: the model config declared BOS/EOS $=1/2$, while
the generation path and the tokenizer reserve policy used id $0$ as EOT/PAD. We resolved
it by reading the active vocabulary rather than choosing a convention. The vocabulary
file has \textbf{65529 lines}, numbered contiguously $1..65529$; there is \emph{no} id
$0$; id $1$ is the NUL byte and id $2$ is the SOH byte. So $1$ and $2$ are ordinary
content tokens, and a config declaring them BOS/EOS would mislabel real bytes and teach
a sequence structure the encoder never inserts. The correct convention is \emph{no BOS}
with EOT $=0$. Ids \texttt{65530}/\texttt{65531}/\texttt{65535} lie outside the content
vocabulary and belong to the tokenizer's special-token layer, which is why the two
layers must not be conflated: \texttt{65531} is absent as a content id and present as
the end-of-text special id.

\paragraph{A third instance, in the same shape.} The two defects above are not the only
ones of their kind in this program, and we record a third because it is the cleanest
demonstration that the class is systematic rather than incidental. A GSM8K evaluation
scored a checkpoint at $4/200$, and the score is not a measurement of arithmetic
ability. The evaluation's prompt asks explicitly for the answer in
\texttt{\textbackslash boxed\{\}}, and classifying all two hundred completions
\textbf{by the scorer's own extraction tiers} gives: the boxed tier fired on
\textbf{zero} items, the last-number fallback on $117$ (of which $4$ were scored
correct), a last-line fallback on $76$ ($0$ correct), and $7$ completions carried no
text at all. So \emph{none} of the four successes came from the tier the prompt
requests; all four came from a last-number rule reading a number out of generative
noise. That the score is noise rather than ability is testable, not merely arguable:
under the null that the extractor draws uniformly from the numbers each completion
happens to emit --- \emph{scoped to the items the fallback tier scored}, since the
null's claim is only that the fallback matched by chance --- the expected count is
$3.60$ against an observed $4$, with $P(X \geq 4) = 0.52$ --- inseparable. (The
scoping is a no-op on this cell: the requested tier fires zero times, so every item
enters the null. It matters for the general form, because a real worked solution
repeats its answer, so an unscoped null's expectation grows with the very format
correct behaviour produces and would eventually flag a model that works.) The completions are degenerate rather than
merely wrong, with $33\%$ spending at least three quarters of their tokens on one
repeated word, a longest identical run of $511$, and extracted values including
\texttt{'.'} and \texttt{'\$'}. The figure is reported nowhere in this paper --- our
descriptive measurements are OpenBookQA and RACE --- and we record it because the
lesson is the same one the decode-side stop set taught: \emph{a number can be
perfectly reproducible and still not be about the thing its label names}. A detail
worth keeping: the first statement of this finding counted whether the
\texttt{\textbackslash boxed} \emph{substring} appeared --- it appears once, in a
completion that emits the token with no braces, which the scorer's brace-balanced
parser therefore rejects. Counting a substring is not counting what the code does,
and that distinction is what separates the two versions of this paragraph.

\paragraph{Why we report this.} The two errors share a shape: a plausible constant,
never checked against the artifact that defines it, in a position where being wrong is
silent. One of them cost three rounds of work in a direction that could not have
succeeded. The paper's later claims inherit the corrected convention, and we record the
correction rather than the repaired number.

\paragraph{The same shape, found in this manuscript.} While preparing
\S\ref{sec:architecture} we checked its description against the code it describes and
found two claims the artifact contradicts; both are corrected in place there. The fusion
gate was written as a function of the two directional outputs --- it is a function of the
shared block input, mixing those outputs convexly --- and the repository's own
architecture note already had it right, so the manuscript was the outlier. The paragraph
on document handling claimed the scan respects document boundaries and that the token
shift is zeroed at each document start; neither holds, and the kernel mode that would
provide per-segment zeroing is never enabled anywhere in the codebase. We record these
for the same reason as the two above, and note one more property they share with them: a
description is a claim like any other, and the isolation sentence had already survived
internal review because it described something obviously desirable rather than something
measured.

\section{The architecture under study}
\label{sec:architecture}

The object is a bidirectional RWKV-7 backbone performing masked diffusion directly over
the token embedding. Each layer holds two distinct RWKV directional modules: a forward
module scanning $x_1,\dots,x_T$ from a zero state, and a backward module scanning the
reversed sequence, also from a zero state, whose outputs are reversed back into token
order. The directions are fused by a learned gate applied to the \emph{shared block
input} $x$ rather than to the two directional outputs:
$\alpha_t=\sigma(W x_t + b)$ and
$o_t=\alpha_t o^{\mathrm f}_t+(1-\alpha_t)o^{\mathrm b}_t$ --- a convex mix, with $W$
zero-initialised against a constant bias so that $\alpha\approx0.98$ at initialisation,
i.e.\ an almost-causal start. Neither direction continues a persistent generation state
across calls; both compute independent, zero-initialised recurrences over the available
window.

\paragraph{What is not isolated.} The recurrence is \emph{not} document-isolated, and we
state this because it bounds what the architecture can be claimed to do. The backbone's
forward signature carries no document boundaries --- no segment lengths and no attention
mask --- and the implementation documents the consequence directly: ``RWKV state simply
flows through pad positions''. A packed row's recurrent state therefore carries across
document boundaries and through padding. Document spans exist in the \emph{data} path
(\texttt{doc\_spans} as \texttt{(doc\_starts, prompt\_ends, doc\_ends)}), where they gate
one optional corruption lane that masks a single packed document's completion; the
default block-wise corruption is gated only by padding, so it can and does mask a
contiguous span across a boundary. Any retention measurement taken on packed rows
therefore includes cross-document state leakage, which is why the registered protocol
measures access modes explicitly rather than assuming isolation.

\paragraph{Depth recycling.} The recycled variant re-applies a \emph{tied} block beyond
the backbone's layer count --- the contract records 32 base plus 16 recycled block passes
per full-canvas invocation, distinct from the outer denoising step count --- so executed
depth rises without parameters rising; only one tanh-bounded scalar gate per extra pass is
added, and it is exactly $0$ at initialisation. The capacity control is a
same-executed-depth \emph{untied} arm, registered as A4 with its unequal parameter count
disclosed rather than matched by construction. A4 is a registered protocol arm and has no
implementation in the released code, so no A4 measurement exists --- the allocation
ledger prices it as a declared derivation from A3 rather than a measurement.

\paragraph{Contract.} The architecture is fixed by a machine-checked contract
(\texttt{DAN/nonlatent\_iclr/architecture\_contract.json}) with 102 passing tests. The
contract records, among its acceptances, the items it does \emph{not} qualify: full-model
FFN token-shift state and tied-loop layer-state alias caches remain
\texttt{NOT\_QUALIFIED} and are explicitly rejected rather than assumed, and
\texttt{whole\_architecture\_ready} remains \texttt{false}. A prefix-cache claim is
therefore unavailable until those modes are wired and qualified, and no result here
depends on them.

\section{Assets, tokenizer, and evaluation}

The evaluation stack is built to be rights-cleared and replayable. RULER
(\texttt{NVIDIA/RULER@c3f5e3b4}, Apache-2.0) and LongBench
(\texttt{THUDM/LongBench@2e00731f}, MIT) are pinned by revision; a 200-task
repository-execution panel is assembled with per-record licences drawn only from
\{MIT, Apache-2.0, BSD-3-Clause\}; in the frozen panel every one of the 200 records is
MIT, and the wider set is the assembly's admission rule rather than a description of the
mix. A deterministic exact-task generator covers associative recall,
overwrite/delayed-query tracking, HMM belief tracking and code-dataflow dependency
queries. The token-length matrix spans 3{,}600 registered units over 720{,}000
instances, with 689 feasible and 31 infeasible cells checked exactly rather than
approximated. Hold-out is by \emph{source group}, not by adjacent chunk, so a reduced
representation cannot be validated on points drawn from a trajectory it already saw.

\paragraph{Recorded limitations.} Three are carried forward unmodified because repairing
them would change what was measured. (i) The repository confirmation split holds only
two tasks (train 184 / validation 14 / confirmation 2), too thin for a repository
endpoint without later panel expansion. (ii) The evaluator is a process-level sandbox
(\texttt{python\_isolated\_mode} + \texttt{posix\_rlimits} + socket stub); this host
denies \texttt{unshare -n}/\texttt{-m} and a writable \texttt{/sys/fs/cgroup}, so it is
\emph{not} an externally managed isolation boundary. (iii) The token matrix ran on a CPU
host rather than in the designated CPU lane.

\section{Registered protocol}
\label{sec:protocol}

Everything in this section is frozen before the confirmatory runs and carries the
failure interpretation agreed in advance. Results appear in \S\ref{sec:notrun}, which is
empty.

The protocol is machine-readable across two files, and they carry different halves of it:
\texttt{DAN/nonlatent\_iclr/preregistration.json} holds the arms, seeds, statistics and
decision rules, with its builder and checks in
\texttt{scale/experiments/nonlatent\_iclr/preregistration.py}; the long-context contract
below --- lengths, loads, cells and instance counts --- is carried by
\texttt{DAN/nonlatent\_iclr/task\_registry.json}, whose computed fields are re-derived by
live replay rather than stored.
It carries \texttt{seal\_state: DRAFT\_UNSEALED} and cannot state otherwise while a
blocker is open: sealing is derived from the inputs, not asserted by a caller, and the
record is written once --- editing it after the fact is detected by digest. Three of its
four sealing blockers remain open (the development-derived H100 deadlines, the frozen
common token budget, and external checkpoint identities); the fourth, the BOS/EOS
generation protocol, closed on 2026-09-15 and is retained in the record as
\texttt{RESOLVED} with its evidence rather than deleted, because a blocker that vanishes
without a trace is indistinguishable from one that was edited away.

\paragraph{The protocol is costed, not aspirational.} The published allocation ledger
measures every arm's throughput at the $4{,}096$-token canvas (A1 $799.1$, A2 $886.7$,
A3 $958.4$, A5 $818.2$ tokens/s; A0 and A4 are declared derivations from A1 and A3, not
measurements) and prices three candidate common budgets. On its per-GPU basis with no
scaling factor, three seeds over six arms cost \textbf{11{,}567} GPU-hours at $2$B
tokens, \textbf{23{,}131} at $4$B, and \textbf{46{,}263} at $8$B. The ledger prices $4$B
for all six arms, so that is the proposal this record carries --- a proposal and not a
freeze, because choosing the budget is a scientific and financial decision with a
$4\times$ range and this record must not make it. The $2.9$B confirmation runs are
\emph{not} covered: the looped large arm reached only $512$-token rungs on a single
$80$\,GiB GPU, so a sharded calibration is required rather than an extrapolation.

\paragraph{Controlled arms.} A0 causal RWKV (autoregressive control); A1 forward-only
denoiser (objective control); A2 bidirectional denoiser (directionality control); A3
bidirectional tied-loop denoiser (the candidate); A4 bidirectional denoiser with untied
extra blocks at equal executed depth (capacity control, unequal parameters disclosed);
A5 forward-only tied-loop (the loop~$\times$~directionality interaction). Seeds 17, 29,
43 for training comparisons; synthetic-data seeds 101--105. Matched-nonpadding-token and
matched-measured-compute readouts are published \emph{separately}, because the two
controls answer different questions and reporting one as the other is the error this
design exists to prevent.

\paragraph{Long-context contract.} Lengths 4{,}096 / 8{,}192 / 16{,}384 / 32{,}768 /
65{,}536 tokens; evidence at 10\%, 50\% and 90\% of the prefix; binding loads
$\{1,8,32,128\}$; distractor loads \{zero, matched random, adversarially similar\}; 200
generated instances per registered cell across five data seeds. Access modes are
full-canvas rescanning, state-only streaming with old tokens inaccessible, and external
retrieval as a separately labelled control --- and a failure of the streaming
implementation does \emph{not} authorise presenting a full-canvas run as
constant-history access. The primary aggregate is an equal-weight macro score over
associative recall, overwrite and dataflow at 16K/32K/64K. OOM and unsupported lengths
remain in the tables as failures, not omitted rows.

\paragraph{Deployment protocol.} H100 measurements are resource-constrained
\emph{proxies} for edge-oriented workloads. Two lanes: independent private-weight
replicas at 4/8/16\,GiB device-resident budgets, and shared-weight serving with
per-session state. Budgets include weights, quantisation metadata, recurrent states,
caches, canvas, logits, activations, workspaces, graphs and framework allocations --- an
allocator-only cap is not a total-memory guarantee. A cell that fails is reported as a
failure. Concurrency grid 1--128, a fixed equal-count 4K/16K/32K/64K mixture, output cap
256 tokens, warm-up 60\,s then at least 300\,s and 1{,}000 completed requests, five
repetitions with randomised arm order. A cell that cannot reach the minimum sample is
labelled \emph{underpowered} rather than given a tail estimate.

\paragraph{Decision rules.} Long context: at least $+5$ absolute points on the primary
macro score with an adjusted confidence interval excluding zero. Serving: at least
$1.5\times$ goodput at a common SLO, or a separately claimed $2\times$ sustainable
session-capacity stretch. Quality equivalence requires the one-sided 95\% lower bound of
candidate-minus-baseline success to be at least $-1$ point. Resampling units are
task/repository clusters, not tokens; 10{,}000 paired hierarchical bootstrap draws with
Holm correction across declared primary comparisons. Seeds are reported individually,
not summarised by a significance label.

\section{Completed runtime qualification}
\label{sec:runtime}

A bounded throughput calibration was completed on $1\times 8$ H100 SXM 80\,GB. It took
six submissions; each failure was diagnosed from its own artifacts, and each fix brought
a CPU regression test. The failures are informative and are reported rather than
elided: a launcher that rejected its own staged root; ranks that found no canvas because
CUDA counters were called without a device; a forward that ran the loop while declaring
zero passes; a per-rung backward warm-up that doubled live activations and OOM'd; a
warm-up sitting outside the OOM guard and discarding earlier rungs; and a caught OOM
followed by further kernels, producing an illegal memory access. The working
configuration --- warm-up inside the OOM guard with the ladder stopping there, allocator
cache released between rungs, a device-bound synchronise, no allocator override ---
passed \textbf{8/8 ranks} with every controlled arm measured at the 4{,}096-token
training canvas.

The resulting allocation ledger prices all six arms, marking A0 and A4 as \emph{declared
derivations} from A1 and A3 rather than measurements. What it does \emph{not} cover is
as important: the 2.9B confirmation runs are not priced, because the looped large arm
OOM'd above 512 tokens on a single 80\,GiB GPU and no sharded calibration has been run.
The ledger states this rather than extrapolating, and it claims no long-context,
goodput, quality or physical-edge result.

\section{A descriptive capability observation}
\label{sec:descriptive}

One set of numbers exists and is worth reporting, provided it is reported for what it
is. On two easy multiple-choice suites --- OpenBookQA and RACE --- the loop s9500
endpoint measures \textbf{56.60} and \textbf{55.40} respectively, against in-lineage
comparators of 54.80/48.00 (b3, s1000), 45.40/45.20 (b3, s8000), 58.40/54.80 (f2,
s14000) and 46.40/31.80 (a released RWKV-7 reference).

Scoped exactly, over the thirteen readings in that summary table: RACE 55.40 is
\textbf{2nd of 13}, behind the released diffusion baseline at 83.60 and ahead of every
other in-house run --- i.e. \textbf{first within this program's own lineage}. OBQA 56.60
is \textbf{4th of 13} and third within the lineage, behind that same released baseline
(84.80) and two f2 runs (60.20, 58.40). We state the ranks because the naturally
stronger sentences are false: ``highest of any model measured'' holds for the lineage
and not for the table, and ``second on OBQA'' holds for neither. Both were asserted
before this line was checked against the table, and both are corrected here --- the
second because the comparison set had been assembled by hand from four models and then
reported as though it were the table.

This is a \emph{descriptive} observation and it is not evidence for the paper's thesis.
It is single-seed; the comparators sit at different training steps and token budgets;
and released checkpoints and in-house runs are not the same-data comparison the
registered protocol requires. Under \S\ref{sec:protocol} it belongs to the
released-versus-controlled result split as description, not as a confirmed effect, and
it moves no endpoint in \S\ref{sec:notrun}.

We include it because reporting only the MMLU curve, which is flatter, would understate
a real and recorded measurement --- and because stating the scope of a favourable-looking
number is the same discipline as \S\ref{sec:correction}. A caveat on its auditability:
the per-example files behind it (\texttt{obqa\_correct/wrong.jsonl}) sit in the
originating program's evidence tree on the cluster, \emph{not} in this repository, so
the number is reproducible from the recorded evaluation command but is not re-derivable
from the published tree alone. They are withheld for the same reason the rest of the
evaluation data is: they contain benchmark items and answers.

\section{Theory: targets, not results}
\label{sec:theory}

The supplied spectral-submanifold material motivates an analysis; it does not prove
retention for a discrete RWKV diffusion model, and treating it as a theorem would be the
error we most want to avoid. The target is a task-specific finite-horizon relation of
the form
\[
E_{k+1}(L) \;\le\; q_k E_k(L) + a_k\,\epsilon_{\mathrm{memory}}(L)
            + \epsilon_{\mathrm{schedule},k} + \epsilon_{\mathrm{model},k}
            + \epsilon_{\mathrm{quant},k},
\]
with every norm, comparator, domain and constant defined, and with the memory term
\emph{derived} from input-dependent recurrent dynamics and a known task statistic rather
than assumed to be small. Under $0\le q_k\le q<1$ and a bounded disturbance the
elementary consequence is $E_K \le q^K E_0 + b(1-q^K)/(1-q)$; that inequality is not the
novelty, and architecture-specific non-vacuous constants with held-out predictions are
what would be. Two constraints are registered in advance: no claim that iterative
refinement recovers information lost through a finite-precision state-only interface,
and no inference of global superiority from local spectral stability. Fitted eigenvalues
inside the unit circle do not justify assuming the full dynamics are contractive.

\section{What is not yet measured}
\label{sec:notrun}

\begin{center}
\begin{tabular}{@{}lll@{}}
\toprule
Endpoint & Registered target & State \\
\midrule
Depth-recycling quality gain & survives matched-compute comparison & \notrun \\
Long-context macro score & $\ge +5$ points, CI excludes zero & \notrun \\
Serving goodput & $\ge 1.5\times$ at common SLO & \notrun \\
Session capacity & $2\times$ stretch, quality non-inferior & \notrun \\
Proxy memory matrix & complete feasibility/precision/context & \notrun \\
Theory & non-vacuous constants, held-out predictions & \notrun \\
Adaptive inference & quality gain at matched total compute & \notrun \\
\bottomrule
\end{tabular}
\end{center}

The confirmatory campaign is not gated on resources. It is gated on a preregistration
that cannot yet be sealed, because sealing requires development-derived H100 SLO
deadlines (which need the deployment measurements), the frozen common token budget, and
resolution of checkpoints that are not present on disk. Until then the protocol above is
a specification, and a specification is not a result.

\paragraph{One result in flight, described as such.} A continuation run
(\texttt{job-49e48c2d}, resumed from step 9500 and still running at its 36-hour
wall-clock cap) is mid-run with a validation masked-CE near 3.81 --- inside the lineage's
3.81--3.83 band and therefore indistinguishable from it. It is \emph{in flight}, it is
not evidence for or against depth recycling, and nothing in this paper depends on it.

\paragraph{What those steps did and did not consume.} This lineage's step counts are
\emph{token-steps}, not data exposure, and the gap is large enough to change what a
token count means --- so the distinction is stated rather than left to the reader.
\texttt{tokens\_seen} is steps $\times$ batch. The data sampler seeds its permutation
from \texttt{(seed, epoch)}, and until 2026-09-17 the trainer never advanced the epoch
and no data position was persisted anywhere: every resume rebuilt the identical
permutation and restarted at row~0, and the iterator was re-created at exhaustion
pointing at the same rows. Measured on this continuation: about $1.43$M unique rows
touched of a $24.39$M-row blend, \textbf{5.9\%}, across six resumes, against a schedule
that claims $5.3$M row-reads. Every ``$N$ tokens'' figure for this lineage should
therefore be read as ``$N$ token-steps over a small fraction of the blend, several
times over''. The defect is fixed in the sampler and the launcher rather than in the
trainer, whose sha256 is pinned by the architecture contract, and the fix is exercised
by every segment boundary of the arm matrix of \S\ref{sec:protocol}.

\section{Reproducibility}

``In flight'' is the honest description of the whole program, so the repository is
organised around making that checkable rather than rhetorical. Every claim above points
at an artifact: a contract with a receipt, a registry whose computed fields are
re-derived by live replay rather than stored, an allocation ledger with planted-violation
refusals, and per-attempt receipts recording the command, environment, input and output
hashes, and the gate verdict. Failed attempts are retained, not overwritten. The code,
contracts and receipts accompany this manuscript; the receipts are hash-bound to the
sources that produced them, and the test suite re-checks those bindings, so the
correspondence between this paper and its evidence is machine-verifiable rather than
asserted.

\paragraph{That correspondence, actually checked.} The sentence above was itself an
assertion until it was run. Forty substantive numbers this manuscript cites are listed
in \texttt{DAN/nonlatent\_iclr/paper\_claims.json} against the artifact that should carry
each one and the \emph{context} it must appear in, and
\texttt{scale/experiments/nonlatent\_iclr/paper\_evidence\_check.py} checks all of them
deterministically, with no model in the loop. The numbers the manuscript uses as
protocol constants rather than as measurements --- the length and load axes, the
bootstrap draw count, the checkpoint step labels --- are declared as such in the same
file, each with a stated reason, so an exemption is reviewable rather than silent. It returns five verdicts, and the
distinctions between them are the point: a value is \emph{verified}, \emph{correctly
rounded} from the artifact's full precision, \emph{derived} by arithmetic over artifacted
inputs, \emph{external} (present only outside the published tree), or \emph{missing}.
Demanding a context rather than accepting a substring is what keeps it honest: on this
tree \texttt{56.60} occurs inside an unrelated token count and \texttt{200} inside a
record identifier, so a bare search would report passes it had not earned. The check is
currently complete --- no value missing, no path missing --- and the only disclosures are
three ledger totals, which are sums the ledger does not state.

\paragraph{A gate must be run on the success case, not only the failure case.} A check
that cannot fail is worthless, but a check that cannot \emph{pass} is worse, because
its failure is indistinguishable from the negative result it destroys: it surfaces
only after the work it was meant to certify. The rule this program now applies is
that every gate is exercised on a case drawn from the real target distribution
\emph{before} the spend it gates, not merely on a case designed to trip it. We record
it because our own evidence checker violated it: it reported a number that was simply
not yet catalogued as a failure of coverage, so a new and perfectly correct
measurement would have read as a defect until someone extended the inventory. The
verdict is therefore three-valued --- the check \emph{found} a problem, or the check
has \emph{not yet been extended} to a number, or it is clean --- and only the first is
a finding.

\paragraph{What each claim is allowed to say.} Existence and support are different
questions, and \texttt{DAN/nonlatent\_iclr/claim\_ceiling.json} answers the second by
declaring every claim against the \emph{kind} of evidence behind it: a real benchmark
reading is capped at the scope actually run, a proved result is capped at the hypotheses
it was proved under, arithmetic over the configuration may not be called a measurement,
and a claim whose experiment has not happened is \emph{blocked} rather than merely
absent. The distinction between blocked and not-applicable is deliberate: an H100 proxy
result can never become an edge-device result, so that claim is out of scope, whereas the
long-context, serving and depth-recycling claims describe work this program intends and
has not yet completed. Three of the eight declared claims are in that blocked state, and
they are the three the confirmatory campaign exists to fill.

The check is only worth its verdict if the verdict is about the manuscript rather than
about the list, and the first version of it was not: it reported completeness while
\S\ref{sec:limitations} cited an endpoint step that no claim carried, because it
audited the inventory and never read the paper. It now extracts every substantive
number from this file and requires each to be either catalogued with a source or
declared structural with a reason, so \emph{complete} means the manuscript is covered
and not merely that the list is.

It also found one real gap, which is why it is worth running. The loop s9500 OpenBookQA
and RACE readings of \S\ref{sec:descriptive} occurred \emph{nowhere} in this repository:
the comparators were recorded, the job that produced the readings was recorded, and the
readings themselves were not. They are now transcribed into
\texttt{DAN/nonlatent\_iclr/descriptive\_readings.json} together with the source path, its
sha256 and its mtime, so the values cited here are checkable here even though the
measurement is not reproducible from this repository alone.

\paragraph{A worked example of the discipline.} When this repository was prepared for
publication, an attempt was made to make host-specific path constants configurable. Ten
source-drift tests failed --- correctly, because those hashes exist to prove the evidence
came from exactly those sources. The change was reverted rather than the hashes
re-derived: re-deriving them would have made the receipts assert a provenance they do not
have. Portability was fixed on the side that carries no provenance claim (the tests)
instead. We report this because it is the same principle as \S\ref{sec:correction}: when
a check and a convenience disagree, the check is the one that is doing work.

\section{Limitations}

Beyond the three recorded asset limitations of \S5: the model has not been shown to
qualify prefix caching, so no constant-history claim is available; the only live
measurements are bounded runtime calibrations rather than capability results; no
checkpoint comparison against external baselines has been run, and any future such
comparison is external validity rather than a same-data causal comparison; the theory
remains a target; and the confirmatory endpoints are all \notrun. A reader should treat
this paper as a corrected foundation plus a frozen protocol, and should not cite it for
an effect size.

\begin{thebibliography}{9}
\bibitem{diffumamba} DiffuMamba: diffusion language modelling with a state-space
backbone. arXiv:2511.15927.
\bibitem{tripletrwkv} Triplet-Block Diffusion RWKV. arXiv:2605.25969.
\bibitem{partialbidir} Partial-bidirectional hybrid block diffusion. arXiv:2607.02805.
\bibitem{pile} L. Gao et al. The Pile: An 800GB dataset of diverse text for language
modelling. arXiv:2101.00027.
\bibitem{ruler} NVIDIA RULER: long-context evaluation, \texttt{NVIDIA/RULER@c3f5e3b4}.
\bibitem{longbench} THUDM LongBench, \texttt{THUDM/LongBench@2e00731f}.
\end{thebibliography}

\end{document}
